\documentclass[10pt,a4paper]{amsart}
\usepackage[margin=19mm]{geometry}
\usepackage{amsmath,amssymb,mathtools}
\usepackage[scr=rsfs,cal=euler]{mathalfa}
\usepackage{xfrac}
\usepackage[unicode,hidelinks,bookmarks=false]{hyperref}
\newcommand{\ie}{i.e.\ }
\newcommand{\cf}{cf.\ }
\newcommand{\resp}{respectively\ }
\newcommand{\Subsection}{\subsection}
\providecommand{\nobreakdash}{\nobreak\hbox{-}}
\setlength{\emergencystretch}{2em}
\multlinegap0pt
\usepackage{amssymb}
\usepackage{enumerate}
\newtheorem{lem}{Lemma}
\newcommand\dbm{d_{\mathrm{BM}}}
\title{Spaces with the maximal projection constant revisited}
\author{Tomasz Kobos}
\date{Source: arXiv:2505.24526v2 (CC BY 4.0)}
\begin{document}
\maketitle

\noindent\emph{Source and licence.} Spaces with the maximal projection constant revisited. Version arXiv:2505.24526v2 (CC BY 4.0), distributed by arXiv under the Creative Commons Attribution 4.0 International licence, http://creativecommons.org/licenses/by/4.0/, as recorded in the arXiv licence metadata for this exact version and shown on the abstract page https://arxiv.org/abs/2505.24526v2. Full licence notice: \texttt{licenses/arxiv-2505.24526-CC-BY-4.0.txt}.\par\smallskip

\noindent\textit{Editorial scope.} Counted unit: the article's lem lemisometric (source0.tex lines 1134--1139). The article's complete original proof of it is delivered with it (source0.tex lines 1141--1180, one \texttt{proof} environment of 40 lines). No mathematics is written, repaired or re-derived: the statement and the proof are the article's own lines, in the article's own notation, and no retained line is substituted or re-flowed.  No further source-local context is needed: every cross-reference of every retained line resolves inside the two delivered spans.  Nothing was omitted from any retained span, so the package carries no omission record.  No retained line contains a citation, so the wrapper carries no bibliography and no bibliography entry.  No retained line contains a figure environment, an image-inclusion command or drawing code, so no image is delivered and the package declares no asset dependency.  Editorial formatting only, in addition: an AMS \texttt{amsart} preamble that restates the article's own declarations and macros used by the retained lines, namely the environments \texttt{lem} on separate counters, together with the standard packages those lines need.  The retained environments print in this document as Lemma 1 in order of appearance, whereas the article prints the same items with the numbering of its own full document.  The complete native source of the article as distributed in the arXiv e-print for this exact version is bundled unchanged as \texttt{sources/179044/source0.tex}.
\begin{lem}
\label{lemisometric}
Let $n \geq 2$ and let $X_0 = (\mathbb{K}^n, \| \cdot\|_0)$, $X_1=(\mathbb{K}^n, \| \cdot\|_1)$ be two normed spaces. For $t \in [0, 1]$ we define a normed space $X_t = (\mathbb{K}^n, \| \cdot \|_t)$ with the norm
$$\|x\|_t = (1-t)\|x\|_0 + t\|x\|_1 \qquad \text{ for } x \in \mathbb{K}^n.$$
If the spaces $X_0$ and $X_1$ are non-isometric, then among the spaces $X_t$ for $t\in[0,1]$ there exists a subset of cardinality continuum consisting of pairwise non-isometric spaces.
\end{lem}

\begin{proof}

We define a mapping $f:[0, 1] \to \mathbb{R}$ as
$$f(t) = \dbm(X_t, X_0).$$
Clearly, $f(0)=1$, $f(t) \geq 1$ for every $t \in [0, 1]$ and $f(1)>1$ by the assumption. We claim that $f$ is a continuous function of $t$. If $0 < t < s < 1$, then by the triangle inequality
$$f(s) = \dbm(X_s, X_0) \leq \dbm(X_s, X_t) \cdot \dbm(X_t, X_0) = \dbm(X_s, X_t) \cdot f(t).$$
Similarly $f(t) \leq \dbm(X_s, X_t) \cdot f(s)$. Moreover for any $x \in \mathbb{K}^n$
$$\frac{t}{s} \left( (1-s)\|x\|_0 + s\|x\|_1  \right ) \leq (1-t)\|x\|_0 + t \|x\|_1 \leq \frac{1-t}{1-s} \left ((1-s)\|x\|_0 + s \|x\|_1 \right ).$$
This yields
$$\dbm(X_s, X_t) \leq \frac{s-ts}{t-ts}$$
and consequently
$$1 \leq \max \left \{ \frac{f(t)}{f(s)}, \frac{f(s)}{f(t)} \right \} \leq \frac{s-ts}{t-ts}.$$
Since $\frac{s-ts}{t-ts} \to 1$ for $s \to t$, this proves continuity of $f$ in every point of the open interval $(0, 1)$.


To establish the continuity in the endpoints, let us pick numbers $\beta > \alpha >0$ such that 
$$\alpha \|x\|_1 \leq \|x\|_0 \leq \beta\|x\|_1,$$
for any $x \in \mathbb{K}^n$. Then for $t \in [0, 1]$ and $x \in \mathbb{K}^n$ we have
\begin{align*}
\left ( 1 + \frac{t}{\beta} - t \right ) \|x\|_0 &=  (1-t) \|x\|_0 + \frac{t}{\beta} \|x\|_0 \leq (1-t) \|x\|_0 + t \|x\|_1 = \|x\|_t \\
 &\leq (1-t) \|x\|_0 + \frac{t}{\alpha} \|x\|_0 = \left ( 1 + \frac{t}{\alpha} - t \right ) \|x\|_0.
\end{align*}
Thus $f(t) \leq \frac{ 1 + \frac{t}{\alpha} - t}{1 + \frac{t}{\beta} - t}$. This expression clearly approaches $1$ as $t \to 0^{+}$, which proves the continuity in $f$ also in $0$. For the continuity in $1$, we note that for every $t \in [0, 1]$ and $x \in \mathbb{K}^n$
\begin{align*}
\left ((1-t)\alpha + t \right ) \|x\|_1  &= (1-t)\alpha \|x\|_1 + t \|x\|_1 \leq (1-t)\|x\|_0 + t\|x\|_1 = \|x\|_t \\
&\leq (1-t) \beta \|x\|_1 + t \|x\|_1 = ((1-t)\beta + t) \|x\|_1.
\end{align*}
Hence
$$\dbm(X_t, X_1) \leq \frac{(1-t)\beta + t}{(1-t)\alpha + t}$$
and from the triangle inequality for $\dbm$
$$\frac{\dbm(X_0, X_1)}{\dbm(X_t, X_1)} \leq \dbm(X_t, X_0) = f(t) \leq \dbm(X_0, X_1) \cdot \dbm(X_t, X_1),$$
so that
$$\dbm(X_0, X_1) \cdot \frac{(1-t)\alpha + t}{(1-t)\beta + t} \leq f(t) \leq \dbm(X_0, X_1) \cdot \frac{(1-t)\beta + t}{(1-t)\alpha + t}.$$
Since 
$$\lim_{t \to 1^{-1}} \frac{(1-t)\beta + t}{(1-t)\alpha + t} = 1, $$
this yields the continuity of $f$ also in $1$.

To conclude the proof, observe that $f(0)=1$ and $f(1)>1$ by the assumption. As $f$ is continuous and nonconstant, the image $f([0, 1])$ contains some nondegenerate interval, and therefore, it contains continuum many different values. Since the Banach--Mazur distance to $X_0$ is an isometric invariant, spaces $X_t$ corresponding to different values of $t$ are non-isometric to each other. Hence there are continuum many pairwise non-isometric spaces among the spaces $X_t$. This concludes the proof.

\end{proof}

\end{document}
